\section{Construction and the two specializations}\label{sec:construction}

We now construct the eigencomplex to which the preceding results apply.
The initial object comes from the unramified characteristic-zero
correspondence.  A degeneration to a separating node produces enhanced
flags on its normalization.  The torus descent of
Proposition~\ref{prop:torus-descent} then gives an ordinary-flag perverse
eigensheaf in characteristic zero.  A second specialization gives the
sheaf in Theorem~\ref{thm:main}.  In both specializations we must place a
nonzero generic stalk in the closure of the particular special-fiber
bundle stack under consideration; this is where uniformization and
proper bounded Hecke spaces enter.

\subsection{A constructible unramified eigencomplex}

Here is the precise consequence of the characteristic-zero correspondence
that we need.  The compactness assertion is essential: the existence of
an arbitrary ind-constructible eigenobject would not suffice for nearby
cycles or for the arguments of Sections~\ref{sec:detection}--\ref{sec:descent}.

\begin{lemma}\label{lem:compact-eigencomplex}
Let $C$ be a smooth projective connected curve over an algebraically
closed field $K$ of characteristic zero, and let $G$ be simple and simply
connected.  Let $\tau$ be a continuous $\Gd$-local system on $C$, defined
over a finite extension of $\mathbb Q_\ell$, with Zariski-dense image.
There is a nonzero locally bounded constructible geometric $E$-adic
complex $F$ on $\Bun_G(C)$ with a coherent Hecke eigenstructure of
eigenvalue $\tau$.  The eigenstructure uses the relative normalization of
Section~\ref{sec:foundations} and includes all finite sets of legs and
their collision maps.
\end{lemma}

\begin{proof}
We use the characteristic-zero part of the Gaitsgory--Raskin theorem,
namely the equivalence
\[
 \operatorname{Shv}_{\mathrm{Nilp}}(\Bun_G(C))
 \simeq
 \IndCoh_{\mathrm{Nilp}}
       \bigl(\Loc^{\mathrm{restr}}_{\Gd}(C)\bigr),
\]
its $\QCoh$-linearity, and the preservation of compact objects by the
inclusion of the automorphic nilpotent category
\cite[Main Theorem~1.3.9(ii), Lemma~1.3.7, and Theorem~1.1.7]{GR}.
Its scope includes geometric $\overline{\mathbb Q}_\ell$-sheaves and
arbitrary algebraically closed characteristic-zero base fields
\cite[Section~2.3]{GR}.  Auxiliary choices in that construction, such as
a theta characteristic, can be made over $K$.

We identify the external Satake labels with our Hecke input--output
convention.  If our label $V$ corresponds externally to $\alpha(V)$
for a symmetric tensor autoequivalence $\alpha$ (for example a Chevalley
relabeling), we use the external spectral point
$\tau\circ\alpha^{-1}$.  Thus, in the convention used below, universal
evaluation at the point denoted $\tau$ is exactly $V\mapsto V_\tau$.

Put $\mathcal S=\Loc^{\mathrm{restr}}_{\Gd}(C)$ and let
$i_\tau:\Spec E\to\mathcal S$ be the point defined by $\tau$.
We first check that the ind-coherent skyscraper
$\delta_\tau=i_{\tau,*}^{\IndCoh}E$ is a nonzero compact object with
nilpotent singular support.  The connected components of the restricted
stack are indexed by semisimple parameters, and a component containing
an irreducible parameter has a unique isomorphism class of geometric
points \cite[Proposition~3.7.2 and Corollary~3.7.5]{AGKRRV}.
Density makes $\tau$ irreducible: a reduction to a proper parabolic
would put its image in that parabolic.  Moreover, $\Gd$ is adjoint, so
\[
 \operatorname{Aut}(\tau)=Z_{\Gd}(\operatorname{im}\tau)
 =Z(\Gd)=1.
\]
The formal tangent complex at this point is
$R\Gamma(C,\ad(\tau))[1]$
\cite[Section~21.2.1]{AGKRRV}.  Its degree $-1$ cohomology vanishes by
the displayed centralizer calculation.  The invariant perfect form on
the characteristic-zero semisimple Lie algebra and Poincar\'e duality
give
\[
 H^2(C,\ad(\tau))
 \simeq H^0(C,\ad(\tau))^\vee(-1)=0.
\]
Thus there are neither infinitesimal automorphisms nor obstructions,
and this one-point formal component is a smooth formal polydisc with
tangent space $H^1(C,\ad(\tau))$.

For completeness, on the completion of a smooth affine space at the
origin, ind-coherent sheaves identify with ind-coherent sheaves on that
space supported at the origin.  The coherent skyscraper is compact in
this category, and extension to the disjoint union of formal components
preserves compactness; this is the formal-completion description used in
\cite[Sections~21.1.7--21.1.10]{AGKRRV}.  On a smooth formal component
the space of obstruction directions is zero.  Hence the spectral
singular support of $\delta_\tau$ is zero, in particular nilpotent.
Here spectral singular support refers to ind-coherent obstruction
directions, rather than to the microlocal support of a constructible
skyscraper.

Let $F$ be the inverse image of $\delta_\tau$ under the equivalence.
It is nonzero and compact in the nilpotent category, hence compact in
the ambient automorphic category by the cited compactness theorem.
Compact automorphic sheaves are bounded constructible after pullback
to each finite-type smooth chart
\cite[Appendix~F, Section~F.2.2]{AGKRRV}.  This proves the required local
finiteness.

It remains to explain why the construction supplies the whole
eigenstructure.  For $A\in\QCoh(\mathcal S)$ the projection formula gives
the natural module identity
\[
 A\otimes\delta_\tau
 \simeq i_{\tau,*}^{\IndCoh}(i_\tau^*A).
\]
For a finite set $I$ and representations $(V_i)_{i\in I}$, apply this
identity to the universal evaluation object with its lisse dependence
on $C^I$.  Its restriction to $\tau$ is
$\boxtimes_{i\in I}(V_i)_\tau$.  The universal evaluation action is the
Hecke action \cite[Main Theorem~14.3.2 and Section~14.3.3]{AGKRRV}.
$\QCoh$-linearity therefore transports the displayed identity to
\[
 \Hecke_{I,(V_i)}(F)
 \simeq F\boxtimes\Bigl(\boxtimes_{i\in I}(V_i)_\tau\Bigr).
\]
These are restrictions of one universal tensor-evaluation system.
The unit, convolution, permutations, and restrictions to collision
diagonals are consequently transported together, with their coherence
relations.  Converting the universal system to our relative Satake
normalization gives the displayed formula without a moving-leg shift,
as fixed in Section~\ref{sec:foundations}.
\end{proof}

\subsection{The separating node and its bundle stack}

Until the mixed-characteristic construction below, the pointed curve
$(X,x)$ is over $\mathbb C$.  Write $U=X\setminus\{x\}$, and fix a
continuous dense parameter $\sigma$ on $U$, represented by $\rho$.
Twisted nodal curves and
their relation to parahoric level already enter the automorphic gluing
construction of Nadler--Yun
\cite[Section~2.3, Proposition~3.1.5, and Lemma~3.2.3]{NYAutomorphicGluing}.
We give the local calculation for the alcove type needed here, including
the gluing torus.  Take copies $(X_1,x_1)$
and $(X_2,x_2)$ of $(X,x)$ and identify $x_1$ with $x_2$ to form the
stable nodal curve $C_0$.  A projective smoothing
\[
 C\longrightarrow\Spec R_1,\qquad R_1=\mathbb C[[s]],
\]
can be chosen with completed local equation $ab=s$ at the node and
smooth connected generic fiber.  Indeed, deformation theory of a
nodal curve has no obstruction in degree two, and its local node
smoothing parameters are independent; algebraizing an ample line
bundle gives the projective family.

Choose compatible roots $s_n$ of $s$, put $R_n=\mathbb C[[s_n]]$ with
$s_n^n=s$, and use the balanced twisted model $C(n)$ over $R_n$.
Its node chart and coarse coordinates are
\begin{equation}\label{eq:construction-balanced-node}
 \left[\Spec R_n[u,v]/(uv-s_n)\big/\mu_n\right],
 \qquad \zeta(u,v)=(\zeta u,\zeta^{-1}v),
 \qquad a=u^n,\quad b=v^n.
\end{equation}
Away from the closed node this is the coarse family after base change.
This description glues in an \emph{\'etale} node chart: changing branch
coordinates multiplies them by units, whose roots can be taken
\'etale locally, and the ambiguities are precisely the indicated
$\mu_n$-action.  Equivalently, one uses the divisor charts of the two
components of the semistable family.  If $n\mid m$, the power maps
$u_n=u_m^{m/n}$, $v_n=v_m^{m/n}$ give compatible maps of twisted models.
Each model is a proper tame Deligne--Mumford stack with projective
coarse space.  The normalization of $C(n)_0$ consists of $X_1$ and
$X_2$ with an $n$th root-stack point at the marking on each.  This is
the balanced-node construction of
\cite[Theorems~1.9--1.10, Remark~1.11, and Proposition~2.2(ii)]{OlssonTwisted}.

Fix $T\subset B\subset G$ and a strictly dominant cocharacter
$\lambda\in X_*(T)$.  Choose $n$ so large that
\begin{equation}\label{eq:construction-alcove}
 0<\langle\alpha,\lambda\rangle<n
 \qquad(\alpha\text{ a positive root}).
\end{equation}
At the node prescribe the conjugacy class of the homomorphism
$\lambda|_{\mu_n}:\mu_n\to G$, using the first branch coordinate $u$
in \eqref{eq:construction-balanced-node}.  Its centralizer is $T$:
no root is trivial on this subgroup by
\eqref{eq:construction-alcove}, and semisimple centralizers in the
simply connected group $G$ are connected.

Let $\mathcal B_n$ be the open substack of $\Bun_G(C(n)/R_n)$ obtained
by removing all the other types in the closed fiber.  This is indeed
open.  The conjugacy types of $\mu_n\to G$ form the finite set
$T[n]/W$, and are locally constant in families on the residual gerbe;
hence the unwanted types form a closed substack of the closed fiber.
The stack $\mathcal B_n$ is algebraic, locally of finite type, and
smooth over $R_n$.  Apply the Hom-stack theorem
\cite[Theorem~1.2]{HallRydh} to the proper flat finitely presented source
$C(n)$ and the target $BG$, whose diagonal is affine.  This gives
algebraicity, local finite presentation, and an affine diagonal.
The obstruction to
deforming a bundle is in $H^2(C(n)_0,\ad P)$, which vanishes: coarse
pushforward is exact for a tame stack
\cite[Lemma~2.3.4]{AbramovichVistoli}, and the coarse space is a curve.
The same calculation applies to every fiber.  Isomorphism spaces are
representable, with the usual affine diagonal for the bundle stack.
In particular, its smooth charts are available for nearby cycles.
Its geometric generic fiber is the unlevelled bundle stack of the
smooth generic curve.

The following calculation identifies its closed fiber, including the
enhancements that will be needed for torus descent.  Set
$U_B=R_u(B)$, let $B^-$ be the opposite Borel containing $T$, write
$U^-=R_u(B^-)$, and put
\[
 \mathcal A_1^+=\Bun_{G,U_B,x_1}(X_1),\qquad
 \mathcal A_2^+=\Bun_{G,R_u(B^-),x_2}(X_2).
\]
An object of either stack is a bundle with a reduction to the specified
unipotent subgroup at the marked point.  Forgetting the enhancement is
a torsor under $T=B/U_B=B^-/R_u(B^-)$.

\begin{lemma}\label{lem:construction-type-flags}
With these choices there is an equivalence
\begin{equation}\label{eq:construction-special-bun}
 (\mathcal B_n)_0
 \simeq [\mathcal A_1^+\times\mathcal A_2^+/T],
\end{equation}
where $T$ changes the two identifications at the gluing gerbe
simultaneously.  Hecke modifications at an unmarked point of $X_i$
pull back to the usual Hecke modifications on the $i$th factor.
\end{lemma}

\begin{proof}
On the first normalized component put $t=u^n$.  Locally on a test
scheme $S$, choose a frame on the root-disc cover in which the inertia
acts by $\lambda(\zeta)$.  Such frames lift from the origin through
all infinitesimal neighborhoods: the frame spaces are smooth and
$\mu_n$-invariants are exact.  A change of frame satisfies
\begin{equation}\label{eq:construction-equivariant-frame}
 h(\zeta u)=\lambda(\zeta)h(u)\lambda(\zeta)^{-1}.
\end{equation}
In particular $h(0)\in T$.  Passing to the invariant punctured-disc
frame replaces $h$ by
$g(t)=\lambda(u)^{-1}h(u)\lambda(u)$.

We check exactly which series $g$ occur.  For a positive root $\alpha$
put $j=\langle\alpha,\lambda\rangle$.  The $\alpha$-root series in
\eqref{eq:construction-equivariant-frame} and its conjugate have the
forms
\[
 u^j a_\alpha(u^n)\longmapsto a_\alpha(t),
 \qquad a_\alpha\in\mathcal O_S[[t]],
\]
whereas the $-\alpha$-root series have the forms
\[
 u^{n-j}b_\alpha(u^n)\longmapsto t b_\alpha(t),
 \qquad b_\alpha\in\mathcal O_S[[t]].
\]
The toral series are simply series in $t$.  These formulas give the
entire group, including its converse.  Namely, $h(0)\in T$ puts $h$
in the formal big cell $U^-TU_B$, in a fixed root order; the displayed
formulas transform that factorization into precisely the big-cell
factorization of
\[
 I_B=\{g(t)\in G(\mathcal O_S[[t]]):g(0)\in B\}.
\]
Conversely every series in $I_B$ has that factorization since its
reduction belongs to $B\subset U^-TU_B$, and the inverse formulas
produce a regular equivariant $h$.  The calculation is valid over
arbitrary test schemes and identifies the group functors.

Gluing to the complement of the point therefore identifies bundles of
this root type with ordinary $B$-level bundles on $X_1$.  Moreover,
evaluation $h\mapsto h(0)$ becomes the torus projection $I_B\to T$.
Its kernel becomes the inverse image of $U_B$ under $I_B\to B$.
Thus an identification of the bundle on the residual gerbe with the
fixed type torsor corresponds exactly to an enhanced flag.

For the second branch the coordinate action is
$v\mapsto\zeta^{-1}v$.  If its positive coordinate generator is denoted
by $\eta=\zeta^{-1}$, the fiber action is $\lambda(\eta^{-1})$.
The preceding calculation therefore uses $-\lambda$ and gives the
opposite Borel $B^-$.  This explains the sign and the choice of
$\mathcal A_2^+$.  Both gerbe frame spaces have automorphism group $T$
when expressed using the common node generator $\zeta$.

A bundle on the nodal twisted curve is a pair of bundles on its
normalization together with an identification on their residual
gerbes.  This is ordinary nodal gluing in the equivariant chart
\eqref{eq:construction-balanced-node}.  Choosing identifications of
both gerbe restrictions with the fixed type torsor makes that gluing
identification automatic, and changing both choices by the same
element of $T$ leaves it unchanged.  Descent gives
\eqref{eq:construction-special-bun}; there is no additional quotient.
If enhancements are instead written as right actions with the second
identification reversed, the same quotient is written with inversion
on the second torus.  We retain the simultaneous-frame convention.

Finally a modification away from the node identifies the bundles near
the gerbe, and hence transports the two frames and the gluing map.
After pullback to the product it acts on exactly the indicated factor.
\end{proof}

At a scheme-theoretic smooth leg of $C(n)$ away from the node,
Beauville--Laszlo gluing \cite[Section~3]{BeauvilleLaszlo}, applied in a
faithful representation and to its reductions of structure group,
gives the usual relative affine Grassmannian description.  Finite jets
suffice for each Schubert bound.  Its bounded
Hecke maps are proper, its exact-position maps to bundle and leg are
smooth, and modifications preserve the prescribed type.  The closed
leg locus contains $U_1\sqcup U_2$.  Consequently this family has all
the geometric properties required in
Theorem~\ref{thm:specialization-detection} and
Proposition~\ref{prop:hecke-specialization}.  The nilpotent cone and its
dimension bound on its special fiber are those for the torus quotient
in Proposition~\ref{prop:cone-dimension}.

\subsection{Gluing the parameter through the node}

We next construct an unramified parameter on the smooth generic curve
whose specialization on $U_1$ is $\sigma$.  This requires a compatible
tower of finite covers, because the given adic parameter need not have
finite monodromy or descend over a finite extension of the trait.

Choose a based orientation-reversing homeomorphism
$f:U_2^{\mathrm{an}}\to U_1^{\mathrm{an}}$, with paths to tangential
base points chosen so that a positive boundary loop $c_2$ maps to
$c_1^{-1}$.  Riemann existence identifies the \'etale fundamental
groups with the profinite completions of these topological groups.
The induced map on completions is continuous.  Thus
$\rho_2=\rho\circ\widehat f_*$ defines a continuous parameter on $U_2$
with the same compact image as $\rho$, and
\begin{equation}\label{eq:construction-boundary-inverse}
 \rho_2(c_2)=\rho(c_1)^{-1}.
\end{equation}
The homeomorphism need not be algebraic: finite-cover Riemann existence
is what supplies the algebraic local systems.

Let $H=\rho(\pi_1^{\mathrm{et}}(U_1))\subset\Gd(L)$, for a finite
coefficient field $L$, and choose a cofinal decreasing sequence of
open normal subgroups $H_r$.  It can be obtained by choosing an
$H$-invariant lattice in a faithful representation and taking
congruence kernels.  Put $Q_r=H/H_r$.  Both punctured components carry
connected $Q_r$-torsors.  Choose root indices $m_r$ with
$n\mid m_r\mid m_{r+1}$ and divisible by the order of the boundary
image in $Q_r$.  The torsors extend to finite \'etale torsors on the
normalization root stacks of $C(m_r)_0$: locally a Kummer root kills
the finite boundary inertia.  This is the root-stack description of
tame covers \cite[Proposition~A.10]{LieblichOlsson}.

At the node the same generator $\zeta$ acts positively on the first
branch and negatively on the second.  Equation
\eqref{eq:construction-boundary-inverse} therefore says that the
restrictions of these torsors to the gluing gerbe have the same
$Q_r$-action.  Fix their identification by reducing one boundary-fiber
identification modulo $H_r$.  This makes all transition maps compatible.
The two torsors glue to a finite \'etale $Q_r$-torsor $Y_{r,0}$ over
$C(m_r)_0$.  On the node chart, both branch covers become constant
across their origins with the same inertia action; gluing the constant
fibers and then taking the equivariant quotient proves the assertion
at the node as well as away from it.

Proper henselian invariance lifts $Y_{r,0}$ uniquely up to its unique
compatible isomorphism to a finite \'etale torsor $Y_r$ over $C(m_r)$
\cite[Theorem~4.5]{Rydh}.  The theorem applies because these stacks are
proper with finite diagonal over the henselian trait.  Its full
faithfulness lifts the group actions, the transition maps after
pullback to a divisible model, and all their relations.  Fix a
compatible geometric generic field
$K=\overline{\mathbb C((s))}$.  All $C(m_r)_K$ are the same smooth
scheme $C_K$.  Choose a base point over a section specializing in
$U_1$, and compatible points above it in the torsor fibers.  These
framings can be lifted from the original tower on $U_1$, since finite
\'etale covers of the strictly henselian section are constant.
The restrictions of the lifted tower give a continuous homomorphism
\begin{equation}\label{eq:construction-generic-parameter}
 \rho_K:\pi_1^{\mathrm{et}}(C_K)\longrightarrow
 \varprojlim_r Q_r=H\subset\Gd(L).
\end{equation}

We verify density rather than assuming that it survives lifting.
Each $Y_{r,0}$ is connected: its first normalization component is
connected, and every component of the other normalization meets the
gluing fiber.  The stack $Y_r$ is proper and flat, with geometrically
reduced fibers.  Its tame coarse space is also flat over the trait,
since taking invariants is exact and preserves flatness.  On the
connected reduced proper special fiber, $H^0(\mathcal O)$ consists
of constants.  Semicontinuity therefore gives
$h^0(Y_{r,K},\mathcal O)=1$, so the geometric generic torsor is
connected.  Hence \eqref{eq:construction-generic-parameter} surjects
onto every $Q_r$.  Its image is compact, thus closed in $H$, and is
therefore all of $H$.  In particular it is Zariski dense in $\Gd$.

For each representation of $\Gd$, enlarge the finite coefficient field
as needed and choose a lattice stable under $H$.  Each finite
reduction factors through some $Q_r$, and hence extends as a lisse
system on the smooth leg locus after the corresponding finite trait
extension.  The special reductions are those of $\sigma$ on $U_1$
and of $\rho_2$ on $U_2$.  Although the models at the node vary with
$m_r$, they agree away from it.  This proves exactly the
lattice-reduction specialization condition of
Proposition~\ref{prop:hecke-specialization}, with its tensor
compatibilities.  It does not require one finite trait extension for
the entire adic system.

\subsection{Nonvanishing at the prescribed type}
\label{subsec:construction-prescribed-type}

Apply Lemma~\ref{lem:compact-eigencomplex} to $C_K$ and $\rho_K$,
obtaining $F$.  Regard its nearby cycles on the fixed model
$\mathcal B_n$.  To prove that they are nonzero, we must meet the
chosen type in the special fiber.  Nonvanishing on some unspecified
component of the full bundle stack would not give this conclusion.

Choose a point of the product in
\eqref{eq:construction-special-bun}.  A smooth chart of $\mathcal B_n$
through its image lifts this point over the henselian trait, and
therefore supplies a reference bundle $P_0$ on $C(n)$ of the prescribed
type.  Choose also a section $y$ of the smooth scheme locus disjoint
from the node.  Since $F$ is nonzero and locally constructible, it has
a nonzero stalk at a $K$-valued bundle $P$ on $C_K$.

The uniformization argument of Drinfeld--Simpson applies here: a
bundle for a semisimple simply connected group on a smooth affine
curve over an algebraically closed field is trivial.  One may use the
precise affine-curve statement
\cite[Theorem~1.1]{BelkaleFakhruddin}, whose fundamental-group-order
hypothesis is automatic for a simply connected group.  Thus $P$ and
$(P_0)_K$ are isomorphic on $C_K\setminus\{y_K\}$.  An isomorphism
exhibits $P$ as a modification of $(P_0)_K$ lying in a finite Schubert
bound.  The corresponding bounded Hecke space over the reference
input and the section $y$ is proper over the trait.  The modification
therefore extends after a finite extension of that trait.  Its output
still has the prescribed node type, since modification at $y$ leaves
the bundle near the node unchanged.

Choose a finite-type smooth affine chart of $\mathcal B_n$ through
this special output.  After a further henselian lift the section of
the bundle stack lifts to that chart.  Its geometric generic point
has the original nonzero stalk.  The closure of the nonzero-stalk
locus on this chart consequently meets the special fiber.  The
contrapositive of Theorem~\ref{thm:specialization-detection}, with the
Hecke geometry and lattice extensions already verified, now gives
\begin{equation}\label{eq:construction-nonzero-node}
                  \Psi F\ne0\quad\text{on }(\mathcal B_n)_0.
\end{equation}

Proposition~\ref{prop:hecke-specialization} gives its coherent
eigenstructure for all legs in $U_1\sqcup U_2$.  Pull back ordinarily
along the smooth surjection
$\mathcal A_1^+\times\mathcal A_2^+\to(\mathcal B_n)_0$.
The pullback is nonzero by \eqref{eq:construction-nonzero-node}.
Choose a geometric point $(a_1,a_2)$ with nonzero stalk and restrict
along $\mathcal A_1^+\times\{a_2\}$.  We obtain a nonzero locally
bounded constructible complex $F_1^+$ on $\mathcal A_1^+$.
Bounded Hecke pushforward for the first curve commutes with both
pullbacks by proper base change and
Lemma~\ref{lem:construction-type-flags}.  Thus $F_1^+$ has eigenvalue
$\sigma$, with all multi-leg compatibilities.  No assertion about
perversity is needed for these ordinary pullbacks.

\begin{proposition}\label{prop:charzero-existence}
Let $(X,x)$ be a smooth projective pointed complex curve of genus at
least two, let $G$ be simple and simply connected, and let $\sigma$ be
a continuous Zariski-dense geometric $E$-adic $\Gd$-parameter on
$X\setminus\{x\}$, defined over a finite coefficient extension.
If its local monodromy is unipotent, there is a nonzero locally
constructible perverse eigensheaf on $\Bun_{G,B,x}(X)$ with eigenvalue
$\sigma$, nilpotent singular support, and the full coherent Hecke
system in our relative normalization.
\end{proposition}

\begin{proof}
The preceding construction gives the nonzero enhanced-flag
eigencomplex $F_1^+$.  Proposition~\ref{prop:torus-descent} applies
because $\sigma$ is dense and its boundary monodromy is unipotent.
It supplies a nonzero locally constructible perverse eigenobject at
ordinary $B$-level.  Its singular support is nilpotent by
Theorem~\ref{thm:nilpotent-detection}.  All steps used geometric adic
sheaves; Betti realization enters only through the previously proved
perversity and torus-descent statements.
\end{proof}

\subsection{Lifting the curve and the parameter from characteristic
\texorpdfstring{$p$}{p}}

We return to the data of Theorem~\ref{thm:main}.  Let $R=W(k)$.
This is a complete strictly henselian discrete valuation ring, and
$\ell$ is invertible in it.  There is a smooth projective pointed lift
\[
 (\mathscr X,\mathfrak x)\longrightarrow\Spec R
 \quad\text{of }(X,x).
\]
Pointed smooth curves have unobstructed deformations, since
\[
                       H^2(X,T_X(-x))=0.
\]
An ample line bundle lifts as well, and formal algebraization produces the displayed projective
family.  Lift the split group and Borel by their split models over
$\mathbb Z$.  Write $\mathscr U=\mathscr X\setminus\mathfrak x$ and
$K_0=\overline{\operatorname{Frac}(R)}$.

We may identify $K_0$ abstractly with $\mathbb C$ for the
characteristic-zero argument.  Indeed, $k$ is countable,
$W(k)$ is a set of countable sequences in $k$, and it contains
$\mathbb Z_p$, so its cardinality is the continuum.  Its fraction
field and an algebraic closure of that field have the same cardinality.
An uncountable algebraically
closed field of characteristic zero has transcendence degree equal
to its cardinality, which proves the assertion.  This identification
concerns the geometric base field; the adic topology and continuity
of the coefficient representation are retained.

\begin{lemma}\label{lem:construction-mixed-parameter}
The given parameter $\sigma$ lifts to a continuous parameter
$\sigma_{K_0}$ on $\mathscr U_{K_0}$ with the same compact image and
the same full boundary inertia homomorphism, after the compatible
identification of prime-to-$p$ roots of unity.  In particular it is
Zariski dense, its local inertia factors through $\mathbb Z_\ell(1)$
and is regular unipotent, and its lattice reductions specialize to
those of $\sigma$ on $U$.
\end{lemma}

\begin{proof}
Use the compact image $H$ and its quotients $Q_r=H/H_r$ as above.
The corresponding $Q_r$-torsor on $U$ is tamely ramified at $x$.
Since its inertia comes from $\mathbb Z_\ell(1)$, its inertia image
has order a power of $\ell$.  Choose compatible indices
$d_r=\ell^{e_r}$, with $d_r\mid d_{r+1}$, that kill these finite
inertia images.  Each torsor extends to a finite \'etale torsor over
the root stack $X(\sqrt[d_r]{x})$.  These indices are invertible in
$R$, so the relative root stacks
$\mathscr X(\sqrt[d_r]{\mathfrak x})$ are smooth proper tame
Deligne--Mumford stacks over $R$.

Apply proper henselian invariance of finite \'etale covers to lift
the torsors, their actions, and all transition maps to these relative
root stacks \cite[Theorem~4.5]{Rydh}.  The root-stack and tame-cover
formulation is also given in
\cite[Proposition~A.10, Theorem~A.11, and Corollaries~A.12--A.13]
{LieblichOlsson}.  Restricting to $\mathscr U$ and then to $K_0$
gives a compatible tower of finite torsors, hence a continuous
homomorphism $\pi_1^{\mathrm{et}}(\mathscr U_{K_0})\to H$.
These statements require the local ramification index to be prime
to $p$; they put no prime-to-$p$ restriction on $|Q_r|$.

Each special torsor on the root stack is connected and geometrically
reduced, since its dense restriction to $U$ is connected.  The proper
flat tame-coarse-space argument used for
\eqref{eq:construction-generic-parameter} gives connected geometric
generic root-stack torsors.  They are smooth connected curves as
stacks; removing the inverse image of the marked gerbe preserves
connectedness.  The homomorphism therefore surjects onto every
$Q_r$, and compactness again gives image exactly $H$.

The gerbe along $\mathfrak x$ records the inertia map at each finite
stage.  Compatible prime-to-$p$ roots of unity identify its special
and generic generators, and full faithfulness in the lifting theorem
preserves the entire action of $\mu_{d_r}$.  Thus the generic inertia
homomorphism factors through $\mu_{\ell^{e_r}}$ at every finite stage
and agrees there with the original one.  In particular every
$\mathbb Z_a(1)$ factor of characteristic-zero inertia for a prime
$a\ne\ell$, including $a=p$, acts trivially.  Passing to the inverse
limit gives the full formula
\[
 \rho_{K_0}(\gamma)
   =\exp\bigl(t_\ell(\gamma)N\bigr)
 \qquad(\gamma\in I_{\mathfrak x,K_0}),
\]
with the same regular nilpotent $N$, up to the allowed change of
generator.  This also accords with the peripheral compatibility of
specialization in \cite[Lemma~3.5]{LieblichOlsson}.

Finally every invariant-lattice reduction in an algebraic
representation factors through some $Q_r$.  The corresponding lifted
torsor is lisse on $\mathscr U$ and has the required special fiber.
The tower respects tensor operations, so these reductions give the
claimed lisse specialization condition for the whole tensor local
system.
\end{proof}

\subsection{Completion of the proof}

\begin{proof}[Proof of Theorem~\ref{thm:main}]\label{proof:main}
Use Lemma~\ref{lem:construction-mixed-parameter} and identify $K_0$
with $\mathbb C$.  Proposition~\ref{prop:charzero-existence} produces
a nonzero locally constructible perverse eigensheaf $M_{K_0}$ on
$\Bun_{G,B,\mathfrak x}(\mathscr X)_{K_0}$ with eigenvalue
$\sigma_{K_0}$.  Set
\[
 \mathscr A=\Bun_{G,B,\mathfrak x}(\mathscr X/R),
 \qquad M=\Psi M_{K_0}\quad\text{on }\mathscr A_k=\mathcal A.
\]
The relative bundle stack is smooth over $R$: the unlevelled bundle
stack is smooth by $H^2$-vanishing, and adding the flag is a smooth
$G/B$-fibration.  By Proposition~\ref{prop:nearby-foundations}, the
geometric nearby cycles used here are locally constructible and
perverse, without taking inertia invariants.  The lattice extensions
in Lemma~\ref{lem:construction-mixed-parameter} and
Proposition~\ref{prop:hecke-specialization} give
\[
 \Hecke_{I,(V_i)}(M)
 \simeq M\boxtimes\Bigl(\boxtimes_{i\in I}(V_i)_\sigma\Bigr)
\]
for every finite set of legs, with the stated naturality, unit,
convolution, permutation, and fusion compatibilities.

It remains to prove that $M\ne0$.  As in
Section~\ref{subsec:construction-prescribed-type}, we extend a bundle
with nonzero generic stalk by a bounded Hecke modification.  Here the
level datum to extend is an ordinary flag, so properness of $G/B$
will complete the extension.  Choose a $K_0$-valued point $(P,\beta)$
where $M_{K_0}$ has nonzero stalk,
a reference $G$-bundle $P_0$ on $\mathscr X$, and a smooth section
$y$ disjoint from $\mathfrak x$.  The existence of $y$ follows by
lifting a point of $X\setminus\{x\}$ over the henselian base.
Off $y_{K_0}$, the bundles $P$ and $(P_0)_{K_0}$ are isomorphic by
the affine-curve triviality used above.  A bounded Hecke space at
$y$, proper over the reference input, extends this modification
after a finite trait extension.  Its output extends the underlying
bundle $P$.  The generic flag $\beta$ then extends by properness of
the associated $G/B$-bundle over the marked section.  Thus the
point with nonzero stalk extends to a section of $\mathscr A$.
Lifting this section in a finite-type smooth affine chart through
its special value shows that the closure of the generic
nonzero-stalk locus meets the special fiber on that chart.

Theorem~\ref{thm:specialization-detection} now proves $M\ne0$.
Its geometric assumptions hold for this smooth pointed family;
the special-fiber cone dimension and transverse-modification
statements are Propositions~\ref{prop:cone-dimension} and
\ref{prop:transverse-hecke}, under precisely the four characteristic
hypotheses of Theorem~\ref{thm:main}.

Finally Theorem~\ref{thm:nilpotent-detection} applied to this locally
constructible eigenobject gives $\SS(M)\subset\Lambda$ on every
smooth chart.  At ordinary $B$-level the cotangent residue lies in
$\Lie R_u(B_\beta)$, so this cone is exactly $\Lambda_{\mathrm{par}}$
from the theorem.  The construction thus has all the required
properties.  All parameter lifts and nearby-cycle operations are
geometric; no Frobenius structure is involved.
\end{proof}
